\chapter{Conclusion}
\label{chap:conclusion}

This chapter synthesizes the main findings, contributions, and limitations of the study, and outlines directions for future research.

\section{Summary of Main Findings}

Using a difference-in-differences framework on 35,264 firm-year observations (2010--2022) and NLP-based digital-transformation measures, this study finds a positive but estimator-sensitive relationship between MSCI inclusion and firm digital-transformation outcomes. The main evidence can be summarized as follows:
\begin{enumerate}
\item TWFE benchmark estimates are positive and statistically strong (0.394 log points in the controlled specification), but cohort-based modern DID diagnostics are materially smaller and less precise, indicating substantial sensitivity to estimator choice under staggered treatment timing;
\item event-study and placebo evidence supports treatment timing plausibility, while alternative outcomes and disclosure-length checks support directional robustness under a hybrid disclosure-substantive measure interpretation;
\item heterogeneity patterns are strongest in high-technology and institutionally developed settings, consistent with absorptive-capacity and institutional-complementarity mechanisms;
\item mechanism evidence is mixed across channels: financing results align more with financing-structure adjustment than broad constraint relaxation, governance proxies suggest strengthened external monitoring conditions with institutional caveats, and direct knowledge-spillover testing remains deferred because key mediators are not yet fully integrated.
\end{enumerate}

Overall, the findings suggest that capital-market liberalization can support digital upgrading, but the average effect is moderate, heterogeneous, and conditional on implementation capacity and institutional context. The study therefore emphasizes calibrated interpretation over universal causal claims.

\section{Contributions}

\begin{enumerate}
\item \textbf{Identification-oriented evidence on digital transformation drivers:} This study provides quasi-experimental evidence that international financial integration is associated with higher firm-level digital transformation intensity in an emerging-market context, while explicitly documenting estimator sensitivity and treatment heterogeneity under staggered adoption.

\item \textbf{Integration of text-as-data with econometric identification:} The study provides evidence that validated NLP-based disclosure measures can be embedded in rigorous DID workflows when accompanied by normalization, placebo testing, and alternative-outcome sensitivity checks.

\item \textbf{Unified mechanism architecture:} Instead of testing channels separately, the study evaluates financing and governance pathways under a common treatment design, generating internally coherent evidence across six directly tested proxy variables.

\item \textbf{Distributional interpretation of policy effects:} By mapping heterogeneity across region, industry, and firm characteristics, the study moves beyond average effects and identifies where liberalization yields the strongest digital dividends and where complementary policy support is most needed.
\end{enumerate}

These contributions are mutually reinforcing: stronger identification supports mechanism interpretation, mechanism evidence improves policy relevance, and transparent measurement improves replicability and external credibility.

\section{Limitations}

Five limitations merit explicit acknowledgment.

First, the text-based outcome captures disclosed digital orientation and intensity rather than direct implementation records. Although normalization and robustness checks mitigate this concern, symbolic disclosure incentives cannot be fully separated from substantive transformation in the current measure.

Second, treatment-effect magnitudes are sensitive to estimator choice in staggered-treatment settings. Modern DID diagnostics indicate that conventional TWFE estimates may overstate average effects when treatment responses are heterogeneous across cohorts and time.

Third, MSCI inclusion is rule-based but not randomly assigned. Treated firms differ systematically from untreated firms on pre-treatment characteristics such as size, liquidity, and governance, and residual bias from unobserved selection cannot be fully ruled out.

Fourth, mechanism measurement remains uneven across channels. Financing and governance proxies are observable in the current panel, but direct knowledge-spillover mediators (e.g., analyst coverage and foreign institutional ownership) are not yet fully harmonized.

Fifth, external validity is bounded by the Chinese institutional context. Financial structure, disclosure regulation, and digital-economy conditions differ across emerging markets, so cross-country generalization should be made cautiously.

\section{Future Work}

Future research can advance along five directions.

First, mechanism identification should extend the baseline by integrating direct knowledge-transfer mediators and estimating the chain MSCI inclusion $\to$ analyst coverage / foreign ownership $\to$ digital transformation.

Second, heterogeneity-robust DID estimators and imputation-based event-study approaches should be implemented as primary estimators when longer post-treatment windows and larger cohort samples become available.

Third, cross-country comparative designs exploiting analogous benchmark-inclusion episodes can clarify which findings are China-specific and which generalize across emerging markets.

Fourth, extending the sample beyond 2022 will allow testing whether estimated effects persist, plateau, or reverse as inclusion factors and global portfolio weights evolve.

Fifth, richer NLP methods such as supervised classification, contextual embeddings, and topic evolution models can improve construct validity by distinguishing strategic rhetoric from operational implementation in corporate digital narratives.

Taken together, the study positions capital market liberalization as a potentially important but conditional contributor to firm-level digital upgrading, with effects that are heterogeneous, estimator-sensitive, and channel-specific in transmission.
